\begin{table}[htbp]
\centering
\caption{Welfare Translation: Averted Deaths and Years of Life Saved}
\label{tab:welfare}
\begin{tabular}{lcc}
\toprule
 & Main design & Dispensary design \\
\midrule
Switch-year ($t=0$) coefficient          & $-2.478$          & $-3.173$ \\
Switch-year counties                     & 160               & 94 \\
Switch-year population                   & 31{,}885{,}361    & 16{,}404{,}374 \\
Averted deaths (95\% CI)                 & 790 (322, 1{,}258) & 520 (136, 905) \\
State-population-weighted life expectancy & 75.8 yrs          & 74.9 yrs \\
National average age at overdose death   & 46.1 yrs          & 46.1 yrs \\
Years of life saved                      & 23{,}462          & 14{,}953 \\
\bottomrule
\end{tabular}
\begin{minipage}{0.92\textwidth}
\vspace{0.3em}
\footnotesize
\textit{Notes:} Averted deaths $= |\hat\beta_{t=0}| \times (\text{population}/100{,}000)$, applied only to the
switch-year event-study coefficient (Tables~\ref{tab:eventstudy_main}--\ref{tab:eventstudy_dispensary}), since that
is the only event-time coefficient distinguishable from zero at conventional levels. Years-of-life-saved
$= \text{averted deaths} \times (\text{state-population-weighted life expectancy} - \text{national average age at
overdose death})$; national average age at death (46.1) is calculated from the 2024 national age-specific overdose
rate distribution. Illustrative of the switch-year point estimate, not a confirmed long-run effect -- Section~4
discusses this caution further, and Section~\ref{sec:cohort} explains why we do not extend this calculation to
the post-switch period.
\end{minipage}
\end{table}
